%\documentclass{clases}
\documentclass[a4paper,12pt, oneside]{book}

\usepackage[skins,minted]{tcolorbox}
\usepackage[utf8]{inputenc}
\usepackage[T1]{fontenc}
\usepackage[spanish, es-tabla]{babel}
\addto\captionsspanish{\renewcommand{\contentsname}{Índice}}
\languageshorthands{spanish}
\usepackage{lmodern} % Usa las fuentes modernas de LaTeX
\usepackage[numbers]{natbib}
%\usepackage{morewrites}
\usepackage{inconsolata}
\usepackage{mdframed}
\usepackage{minted}
\usepackage{bm}
\usepackage{listingsutf8}
\lstset{
  inputencoding=utf8,
  language=Python,
  % … tus otras opciones …
}
\usepackage{subcaption}
\usepackage{amsmath}
\usepackage{amssymb}
\usepackage{graphicx}
\usepackage{hyperref}
\usepackage{longtable}
\usepackage{tabularx}
\usepackage{threeparttable}
\usepackage{array}    % Necesario para centrar horizontal y verticalmente
\usepackage{xcolor}
\usepackage{pdfpages}
%\usepackage{color}
\usepackage{fancyhdr}
\usepackage{menukeys}
\usepackage{appendix}
\usepackage{fontawesome}
\usepackage{comment}
\usepackage{caption}
\usepackage{setspace}
\usepackage[explicit]{titlesec}
\usepackage[a4paper]{geometry}
\geometry{top=3cm, bottom=3cm, left=4cm, right=2cm}

\definecolor{green1}{HTML}{1dae28}
\definecolor{green2}{HTML}{afd095}
\definecolor{lightgray}{gray}{0.9}
\definecolor{orange}{RGB}{18,84,183}
\definecolor{titulo}{gray}{0.75}
\definecolor{gray97}{gray}{.97}
\definecolor{gray75}{gray}{.75}
\definecolor{gray45}{gray}{.45}
\definecolor{advertencia}{RGB}{255,178,102}
\definecolor{colorturqueza}{RGB}{178,223,238}
\definecolor{mintedbackground}{rgb}{0.95,0.95,0.95}
\definecolor{lbcolor}{rgb}{0.95,0.95,0.95}
\definecolor{mintedframe}{rgb}{0.0,0.0,0.0}
\definecolor{codebg}{rgb}{0.96,0.96,0.96}
\definecolor{colorurls}{RGB}{107,17,17}
\definecolor{colorsql}{RGB}{255,245,245}
\definecolor{colorreferences}{RGB}{48,134,3}
\definecolor{titulo}{gray}{0.65}			%------ color para fondo del titulo de tablas.

\hypersetup{
	%bookmarks=true,         % show bookmarks bar?
	unicode=true,          % non-Latin characters in Acrobat’s bookmarks
	pdftoolbar=true,        % show Acrobat’s toolbar?
	pdfmenubar=true,        % show Acrobat’s menu?
	pdffitwindow=false,     % window fit to page when opened
	pdfstartview={FitH},    % fits the width of the page to the window
	pdftitle={Reporte final de Robótica},    % title
%	pdfauthor={},     % author
	pdfsubject={Reporte final de Robótica},   % subject of the document
	%pdfcreator={pdfTeX 3.14159265-2.6-1.40.16 (TeX Live 2016/dev)},   % creator of the document
	%pdfproducer={Panel HJ 2017}, % producer of the document
	pdfkeywords={Manipulador industrial} {Robotica} {ros} {gazebo}, % list of keywords
	%pdfnewwindow=true,      % links in new PDF window
	colorlinks=true,       % false: boxed links; true: colored links
	linkcolor=black,          % color of internal links (change box color with linkbordercolor)
	citecolor=colorreferences,        % color of links to bibliography
	filecolor=magenta,      % color of file links
	urlcolor=blue,           % color of external links
	linkbordercolor={0 0 0}
}

\lstset{
	frame=Ltb,
	tabsize=2,
	framerule=0pt,
	aboveskip=0.5cm,
	framextopmargin=0pt,
	framexbottommargin=0pt,
	framexleftmargin=0.4cm,
	framesep=0pt,
	rulesep=.0pt,
	backgroundcolor=\color{gray97},
	rulesepcolor=\color{blue},
	%
	stringstyle=\ttfamily,
	showstringspaces = false,
	basicstyle=\small\ttfamily,
	commentstyle=\color{gray45},
	keywordstyle=\bfseries,
	%
	numbers=none,
	numbersep=15pt,
	numberstyle=\tiny,
	numberfirstline = false,
	breaklines=true
}

\setminted[bash]{
	bgcolor=mintedbackground,
	fontfamily=tt,
	linenos=true,
	numberblanklines=true,
	numbersep=11pt,
	numbersep=2pt,
	gobble=0,
	frame=leftline,
	framesep=2mm,
	funcnamehighlighting=false,
	tabsize=4,
	obeytabs=false,
	samepage=false,
	showspaces=false,
	showtabs =false,
	texcl=false,
	baselinestretch=1.2,
	fontsize=\footnotesize,
	breaklines=true,
	breaksymbolleft=\ 
}

\lstdefinestyle{consola}{
	basicstyle=\footnotesize\bf\ttfamily,
	backgroundcolor=\color{gray75},
}	
\definecolor{gray}{rgb}{0.4,0.4,0.4}
\definecolor{darkblue}{rgb}{0.0,0.0,0.6}
\definecolor{cyan}{rgb}{0.0,0.6,0.6}
\lstset{language=XML}

\lstdefinelanguage{XML}{
	morestring=[b]",
	tabsize=2,
	breaklines=true,
	morestring=[s]{>}{<},
	morecomment=[s]{<?}{?>},
	stringstyle=\color{black},
	identifierstyle=\color{darkblue},
	keywordstyle=\color{cyan},
	numbers=left,
	morekeywords={xmlns,version,type}% list your attributes here
}

\lstdefinestyle{C}{language=C}
\lstdefinestyle{XML}{language=XML}

\DeclareMathOperator{\diag}{diag}
\lstset{literate=
	{á}{{\'a}}1 {é}{{\'e}}1 {í}{{\'i}}1 {ó}{{\'o}}1 {ú}{{\'u}}1
	{Á}{{\'A}}1 {É}{{\'E}}1 {Í}{{\'I}}1 {Ó}{{\'O}}1 {Ú}{{\'U}}1
	{à}{{\`a}}1 {è}{{\`e}}1 {ì}{{\`i}}1 {ò}{{\`o}}1 {ù}{{\`u}}1
	{À}{{\`A}}1 {È}{{\'E}}1 {Ì}{{\`I}}1 {Ò}{{\`O}}1 {Ù}{{\`U}}1
	{ä}{{\"a}}1 {ë}{{\"e}}1 {ï}{{\"i}}1 {ö}{{\"o}}1 {ü}{{\"u}}1
	{Ä}{{\"A}}1 {Ë}{{\"E}}1 {Ï}{{\"I}}1 {Ö}{{\"O}}1 {Ü}{{\"U}}1
	{â}{{\^a}}1 {ê}{{\^e}}1 {î}{{\^i}}1 {ô}{{\^o}}1 {û}{{\^u}}1
	{Â}{{\^A}}1 {Ê}{{\^E}}1 {Î}{{\^I}}1 {Ô}{{\^O}}1 {Û}{{\^U}}1
	{œ}{{\oe}}1 {Œ}{{\OE}}1 {æ}{{\ae}}1 {Æ}{{\AE}}1 {ß}{{\ss}}1
	{ç}{{\c c}}1 {Ç}{{\c C}}1 {ø}{{\o}}1 {å}{{\r a}}1 {Å}{{\r A}}1
	{€}{{\EUR}}1 {£}{{\pounds}}1 {'}{{\textquotesingle}}1 {Ñ}{{\~N}}1
	{ñ}{{\~n}}1
}

\newtcblisting{terminal}[2][]{
	listing engine=minted,
	listing only,
	#1,
	title=#2,
	minted language=bash,
	colback=mintedbackground,
	top=0mm,
	bottom=0mm
}

\newtcblisting{consolestyle}[2][]{enhanced, listing engine=minted, 
	listing only,#1, title=#2, minted language=bash, 
	coltitle=mintedbackground!35!black, 
	fonttitle=\ttfamily\footnotesize,
	sharp corners, top=0mm, bottom=0mm,
	title code={\path[draw=mintedframe, dashed, fill=mintedbackground](title.south west)--(title.south east);},
	frame code={\path[draw=mintedframe, fill=mintedbackground](frame.south west) rectangle (frame.north east);}
}


\graphicspath{{img/}{tablas/}{portada/}}  % Las imágenes se buscarán en la carpeta "img"

\newenvironment{doble}
{\doublespacing
}

%\newcounter{comando}[section]
%\newenvironment{comando}[1][]{\refstepcounter{comando}\par\medskip
	%	\noindent \textbf{Comando~\thecomando. #1} \rmfamily}{\medskip}
%\begin{terminal}{#1}

%\end{terminal}
%}{\medskip}

\graphicspath{ {img/} }

\begin{document}

	\newfloat{Comando}{h}{lot}[chapter]
	
	\renewcommand{\tablename}{Tabla}
	\renewcommand{\listtablename}{Índice de tablas}
	\renewcommand\listoflistingscaption{Índice de códigos}
	\renewcommand\listingscaption{Código}
	\newcommand{\code}[1]{\colorbox{lightgray}{\texttt{#1}}}
	\renewcommand*{\lstlistlistingname}{Índice de fragmentos de código}

	\renewcommand\lstlistingname{Código}
	\renewcommand{\appendixname}{Anexo}
	\renewcommand{\appendixtocname}{Anexos}
	%\renewcommand{\appendixpagename}{Anexo}
	\def\CC{{C\nolinebreak[4]\hspace{-.05em}\raisebox{.4ex}{\tiny\bf ++ \space}}}
	\def\Cc{{C\nolinebreak[4]\hspace{-.05em}\raisebox{.4ex}{\tiny\bf ++}}}
	\newcommand{\ffolder}[1]{\menu{\faFolderO \space #1}}
	\newcommand{\ffile}[1]{\menu{\faFileO \space #1}}
	\newcommand{\folder}{\faFolderO \space}
	\newcommand{\file}{\faFileO \space}
	\newcommand{\world}{\faGlobe \space}
	\newcommand{\wworld}[1]{\menu{\faGlobe \space #1}}
	\renewcommand\labelitemi{$\bullet$}
	\newcommand{\SB}{\{B\}} % Define el sistema del cuerpo
	\newcommand{\SI}{\{I\}} % Define el sistema inercial
	
	\newcounter{actividad} % Define un contador llamado "actividad"
	
	% Aquí se encuentra el archivo con la portada
	\begin{titlepage}
	\centering
	%-------------------------------------------
	% Logos en una tabla: izquierda, centro y derecha
	\begin{tabular}{@{}p{0.3\textwidth} p{0.3\textwidth} p{0.3\textwidth}@{}}
		\includegraphics[height=2cm]{tecnm} & 
		\centering \includegraphics[height=1.5cm]{SEP} & 
		\raggedleft \includegraphics[height=2cm]{ith.jpg} \\
	\end{tabular}
	
	\vspace{2em}
	
	\noindent
	%-------------------------------------------
	%	Información institucional y académica (esquina superior izquierda)
	\begin{minipage}[t]{0.48\textwidth}
		\raggedright
		\small \textbf{%
			Instituto Tecnológico de Hermosillo\\
			Materia: Robótica\\
			Profesor: Medina Gil Lamadrid, Jesús Iván%
		}
	\end{minipage}%
	\hfill
	%	fecha actual (esquina superior derecha), en letras pequeñas y en negrita.
	\begin{minipage}[t]{0.48\textwidth}
		\raggedleft
		\small \textbf{\today}
	\end{minipage}
	
	\vspace{2em}
	
	%-----------------------------------------
	% Unidad y Título de la tarea en letras grandes y en negrita
	{\large \textbf{Unidad 1: Morfología del robot}}\\
	{\Huge \textbf{Investigar sobre diferentes tipos de sensores}}
		
	\vspace{1em}
	
	%---------------------------------------
	% Tabla con la información del equipo
	%---------------------------------------
	% Encabezado del equipo
	\begin{center}
		{\Large \textbf{Equipo 4}}
	\end{center}
	
	\vspace{1em}
	
	% Tabla de integrantes:
	% Cada fila contiene: foto (columna izquierda) y datos del integrante (columna derecha)
	\begin{center}
		\begin{tabular}{c c}
			\begin{tabular}{c}
				\includegraphics[height=3cm]{perfil1.jpg} \\
				\textbf{Fuentes Ochoa},\\ Aislinn Alicia \\ \texttt{l21330583@hermosillo.tecnm.mx} \\ Teléfono: 6371147080
			\end{tabular} &
			\begin{tabular}{c}
				\includegraphics[height=3cm]{perfil2.jpg} \\
				\textbf{Gonzalez Cueto,}\\ Alejandra Abigail \\ \texttt{l21330591@hermosillo.tecnm.mx} \\ Teléfono: 6221223887
			\end{tabular} \\ \vspace{2em}
			\begin{tabular}{c}
				\includegraphics[height=3cm]{perfil3.jpg} \\
				\textbf{Ceballos Portillo,}\\ Patsy \\ \texttt{l21330551@hermosillo.tecnm-mx} \\ Teléfono: 6622968916
			\end{tabular} &
			\begin{tabular}{c}
				\includegraphics[height=3cm]{perfil4.jpg} \\
				\textbf{Peña Encinas}\\ Ana Lourdes \\ \texttt{l21331075@hermosillo.tecnm.mx} \\ Teléfono: (6621281812)
			\end{tabular}
		\end{tabular}
	\end{center}

\end{titlepage}
	
	\spacing{1.5}
	\frontmatter
	%\chapter{Dedicatoria}
	%\input{capitulos/01_dedicatoria}
	
	%\chapter{Declaración}
	%\input{capitulos/02_declaracion}
	\titleformat{\chapter}
	{\bfseries\huge}
	{}
	{0pt}
	{~\raisebox{-1.5pt}{}
	\\\filleft #1 \\\vspace{.25cm}\titlerule[1.5pt]}
	
	% ---------------------------------------
	% índices
	\newpage\addcontentsline{toc}{chapter}{Índice general}\tableofcontents
	
	%\listof{Comando}{Índice de Comandos}\addcontentsline{toc}{chapter}{Índice de comandos}

	\newpage\addcontentsline{toc}{chapter}{Índice de figuras}\listoffigures
	
	%\newpage\addcontentsline{toc}{chapter}{Índice de tablas}\listoftables
	
	%\newpage\addcontentsline{toc}{chapter}{Índice de fragmentos de código}\lstlistoflistings
	
	\hypersetup{
		linkcolor=red
	}
	
	\mainmatter
	\pagestyle{fancy}
	\fancyhead{}
	\fancyhead[L]{\leftmark}
	\fancyhead[R]{}
	\fancyfoot[L]{\parbox[l]{\textwidth-1cm}{\rightmark}}
	\fancyfoot[C]{}
	\fancyfoot[R]{\thepage}
	\renewcommand{\footrulewidth}{0.5pt}
	%\fancyfoot[RO,LE]{Diseño de modelo para simulación 3D de VANT tipo cuadricóptero}
	
	\titleformat{\chapter}
	{\bfseries\huge}
	{}
	{0pt}
	{\titlerule[3pt]~\raisebox{-1.5pt}{\sc{\chaptername}~\thechapter}~\titlerule[3pt]%
		\\\vspace{.05cm}\titlerule\\\filcenter #1 \\\vspace{.25cm}\titlerule}
	
	\chapter{Introducción} \label{chap:planteamiento-del-problema}
	\chapter{Marco Teórico} 
\label{chap:marco_teorico}

La robótica es una rama de la ingeniería que estudia el diseño, construcción, operación y aplicación de robots, así como los sistemas computacionales que los controlan. Un robot es una máquina reprogramable capaz de realizar una serie de acciones automáticamente o por control externo, con cierto grado de autonomía. El estudio de los robots implica diversas áreas del conocimiento, entre ellas la mecánica, la electrónica, la programación y la teoría del control.

La cinemática se encarga del estudio del movimiento de los cuerpos sin considerar las fuerzas que los causan. En robótica, se distingue entre:

\textbf{Cinemática directa:} determina la posición y orientación del efector final en función de los valores articulares.

\textbf{Cinemática inversa:} calcula los valores articulares necesarios para alcanzar una posición deseada del efector.


El análisis cinemático es esencial para programar tareas de posicionamiento precisas en robots manipuladores.

La dinámica estudia el comportamiento del robot considerando las fuerzas involucradas en su movimiento. Permite calcular los torques o fuerzas necesarias para lograr un cierto movimiento bajo condiciones dinámicas. Este análisis es fundamental para el diseño del sistema de control y para garantizar que el robot pueda ejecutar tareas de manera estable y segura.
Las propiedades dinámicas de un robot incluyen:

\textbf{Inercia:} resistencia de las masas del robot al cambio de movimiento.

\textbf{Centrífugas y coriolis:} fuerzas que aparecen en sistemas en movimiento.

\textbf{Par gravitacional:} el efecto del peso del robot sobre sus articulaciones.

\cite{fundamentosRobotica}
\cite{introduccionRobotica}





	
	\part{Marco teórico}
		\chapter{Revisión del estado del arte} \label{chap:estado del arte}
		\input{capitulos/estado_del_arte}

		\chapter{Dinámica modelada por ángulos de Euler} \label{chap:CMD}
		\input{capitulos/Control_modos_deslizantes}
		
		\chapter{Dinámica modelada por cuaterniones} \label{chap:cuaterniones}
		\input{capitulos/cuaterniones}
		
		\chapter{Estimación de los estados} \label{chap:estimacion}
		\input{capitulos/estimacion_estados}
		
	\part{Desarrollo}
%		\chapter{Introducción a la propuesta de solución} \label{chap:Introducción a la propuesta de solución}
%		\input{capitulos/Desarrollo_Introduccion}
		
		\chapter{Metodología} \label{chap:Metodologia}
		\input{capitulos/Metodologia}
		
		\chapter{Materiales y equipo} \label{chap:Materiales}
		\input{capitulos/Materiales}
		
		\chapter{Avances y resultados} \label{chap:Avances y resultados}
		\input{capitulos/Avances_resultados}
		
%		\chapter{Propuesta de solución}
%		\chapter{Metodología}
%		\chapter{Descripción de las actividades a realizar}
%		\chapter{Cronograma de actividades}
%		\chapter{Presupuesto requerido (equipo y material de laboratorio)}
%		\chapter{Pruebas realizadas}
%		\input{capitulos/Desarrollo}
		
	\part{Resultados, conclusiones y recomendaciones}
		\chapter{Conclusiones y recomendaciones} \label{chap:conclusiones}
		\input{capitulos/Conclusiones}
		
%		\chapter{Competencias desarrolladas y/o aplicadas}
%		\input{capitulos/competencias-desarrolladas}
%	
%	\part{Desarrollo}
%		\chapter{Procedimiento y descripción de las actividades realizadas}
%		\input{capitulos/procedimiento-y-descripcion-de-las-actividades-realizadas}
%		\chapter{Preparación del entorno}\label{chap:preparacion-del-entorno}
%		\input{capitulos/preparacion_del_entorno}
%		\chapter{Proyectos similares}\label{chap:proyectos-similares}
%		\input{capitulos/proyectos-similares}
%		
%	\part{Resultados, conclusiones y recomendaciones}
%		\chapter{Análisis de resultados}
%		\input{capitulos/14_analisis_de_resultados}
%		
%		\chapter{Conclusiones y recomendaciones}\label{chap:conclusiones}
%		\input{capitulos/21_conclusiones_y_recomendaciones}
%		
%		\chapter{Competencias desarrolladas y/o aplicadas}
%		\input{capitulos/competencias-desarrolladas}

	%\printglossary[title=Glosario,toctitle=Glosario]
	%linea que anexa la palbra Glosario al índice
	%\addcontentsline{toc}{chapter}{Glosario}
	
	%\printglossary[type=\acronymtype]
	%linea que anexa la palbra Acrónimo al índice
	%\addcontentsline{toc}{chapter}{Acrónimos}

%	\part*{Bibliografía y Firmas}
	\addcontentsline{toc}{part}{Bibliografía y Firmas}
	%\newpage
	\titleformat{\chapter}
	{\bfseries\huge}
	{}
	{0pt}
	{\titlerule[3pt]~\raisebox{-1.5pt}{\sc{	\appendixname}~\thechapter}~\titlerule[3pt]%
		\\\vspace{.05cm}\titlerule\\\filcenter #1 \\\vspace{.25cm}\titlerule}
	
	
	
	%\chapter{\texttt{fly\_bot}}
	%\input{capitulos/fly-bot}
	
%	\part{Manual de uso de Hector quadrotor}\appendix\addappheadtotoc\label{part.manual}
%		\chapter{Descripción y simulación del modelo de Hector quadrotor}\label{chap:hector-traduccion}
%		\input{capitulos/hector-traduccion}
%		
%		\chapter{\texttt{hector\_quadrotor}}\label{sec.quadrotor}
%		\input{capitulos/A1_manual}
%		
%		\chapter{\texttt{hector\_gazebo}}\label{sec.gazebo}
%		\input{capitulos/hector-gazebo}
%		
%		\chapter{\texttt{hector\_models}}\label{sec.models}
%		\input{capitulos/hector-models}
		
		%\chapter{\texttt{hector\_localization}}\label{sec.localization}
		%\input{capitulos/hector-localization}
		
		%\chapter{\texttt{hector\_slam}}\label{sec.hector.slam}
		%\input{capitulos/hector-slam}
	%\input{capitulos/25_anexos}
	%\chapter{Símbolos}
	%\listofsymbols\chapter
	
	\titleformat{\chapter}
	{\bfseries\huge}
	{}
	{0pt}
	{~\raisebox{-1.5pt}{}
		\\\vspace{.05cm}\titlerule\\\filcenter #1 \\\vspace{.25cm}\titlerule}
	%{\titlerule\\\vspace{.25cm}\filcenter #1 \\\vspace{.25cm}\titlerule}
	\bibliographystyle{IEEEtranN}
	\newpage\label{bibliografia}
	\addcontentsline{toc}{chapter}{Bibliografía}\bibliography{bibliografia}
	\pagestyle{plain}
	\newpage\label{Firmas}
	\addcontentsline{toc}{chapter}{Firmas}
	\includepdf[pages=1, fitpaper=true]{portada/Firmas}
	\spacing{1.5}
\end{document}


